\color{unidarkgreen}{\subsection{Дифференциальная защита шин}}
\color{black}

\begin{enumerate}[label=\arabic{section}.\arabic{subsection}.\arabic{enumi}, labelsep=4pt, leftmargin=0pt, itemindent=57pt, itemsep=0pt, parsep=5pt] % если что можно крутить itemsep

%%%%%%%%%%%%%%%%%%%%%%%%%%%%%%%%%%%%%%%%%%%%%%%%%% ОБШИЕ СВЕДЕНИЯ %%%%%%%%%%%%%%%%%%%%%%%%%%%%%%%%%%%%%%%%%%%%%%%%%%
\item Общие сведения

\begin{enumerate}[label=\arabic{section}.\arabic{subsection}.\arabic{enumi}.\arabic*, labelsep=4pt, leftmargin=0em, itemindent=65pt, parsep=0pt]

\item
ДЗШ является быстродействующей защитой с абсолютной селективностью и выполняет функцию основной токовой защиты шин от всех видов КЗ.

\item
Данное устройство ДЗШ имеет пофазное исполнение, т.е. каждая фаза обрабатывается одним терминалом независимо от двух других.
Для селективной защиты одной фазы устройство содержит три дифференциальных органа: пусковой орган (ПО) для защиты обеих СШ, избирательные органы для защиты первой (ИО1) и второй (ИО2) СШ.
Ввод и вывод данной функции для всех дифференциальных органов осуществляется либо посредством одного программного ключа \textbf{«ДЗШ»}, либо с использованием внешнего сигнала \textit{«Вывод ДЗШ»}.

\item
При рабочей фиксации присоединений сигнал на отключение поврежденной СШ действует только при срабатывании пускового органа и соответствующего избирательного органа.
В режиме нарушенной фиксации избирательные органы выводятся из работы и отключение обеих СШ производится только от пускового органа.

\item
Пусковой орган подключен на токи всех присоединений обеих СШ, за исключением токов ШСВ.
Избирательные органы подключены на токи всех присоединений соответствующей СШ, включая токи ШСВ.
Логика фиксации присоединений приведена в разделе «Логика управления топологией».

\item
В устройстве предусмотрена настройка полярности ТТ присоединений.
По умолчанию измерение токов выполняется в направлении к защищаемому объекту, то есть все ТТ считаются включёнными встречно друг к другу.
Если для какого-либо присоединения фактическая полярность не совпадает с принятой по умолчанию, пользователь может скорректировать её двумя способами: либо изменением подключения вторичных токовых цепей, либо установкой программного ключа \textbf{«Инверсия Bn»} в положение \textbf{«Введено»}.

\small
\begin{longtable}{|>{\centering\arraybackslash}m{6.4cm}|>{\centering\arraybackslash}m{4cm}|>{\centering\arraybackslash}m{2.8cm}|>{\centering\arraybackslash}m{2.8cm}|}
% Первая страница: полный заголовок
\caption{Перечень параметров ТТ и ТН\hfill\vspace{-0.5\baselineskip}}\label{desc:tbl1}\\
\hline
\rowcolor{gray!30}
Наименование уставки & Диапазон & Шаг & Значение по умолчанию \\
\hline
\endfirsthead

% Промежуточные страницы: «Продолжение таблицы»
\caption*{\hspace{3pt}\emph{Продолжение таблицы \ref{desc:tbl1}}\hfill\ }\vspace{-0.5\baselineskip} \\
\hline
\rowcolor{gray!30}
Наименование уставки & Диапазон & Шаг & Значение по умолчанию \\
\hline
\endhead

\centering\arraybackslash Номинальный первичный ток ТТ В1, А  & \centering\arraybackslash 1…12000 & \centering\arraybackslash 1 & \centering\arraybackslash 1000 \\
\hline
\centering\arraybackslash Номинальный вторичный ток ТТ В1, А  & \centering\arraybackslash 1…12000 & \centering\arraybackslash 1 & \centering\arraybackslash 1000 \\
\hline
\centering\arraybackslash Инверсия В1  & \centering Введено \\ Выведено & \centering --- & \centering\arraybackslash Выведено \\
\hline
\centering\arraybackslash …  & \centering\arraybackslash … & \centering\arraybackslash … & \centering\arraybackslash … \\
\hline
\centering\arraybackslash Номинальный первичный ток ТТ В12, А  & \centering\arraybackslash 1…12000 & \centering\arraybackslash 1 & \centering\arraybackslash 1000 \\
\hline
\centering\arraybackslash Номинальный вторичный ток ТТ В12, А  & \centering\arraybackslash 1…12000 & \centering\arraybackslash 1 & \centering\arraybackslash 1000 \\
\hline
\centering\arraybackslash Инверсия В12  & \centering Введено \\ Выведено & \centering --- & \centering\arraybackslash Выведено \\
\hline
\centering\arraybackslash Номинальное первичное линейное напряжение ТН 1С, кВ  & \centering\arraybackslash 0,22…1150,00 & \centering\arraybackslash 0,01 & \centering\arraybackslash 110 \\
\hline
\centering\arraybackslash Номинальное вторичное линейное напряжение ТН 1С, кВ  & \centering\arraybackslash 100,00…400,00 & \centering\arraybackslash 0,01 & \centering\arraybackslash 100 \\
\hline
\centering\arraybackslash Номинальное первичное линейное напряжение ТН 2С, кВ  & \centering\arraybackslash 0,22…1150,00 & \centering\arraybackslash 0,01 & \centering\arraybackslash 110 \\
\hline
\centering\arraybackslash Номинальное вторичное линейное напряжение ТН 2С, кВ  & \centering\arraybackslash 100,00…400,00 & \centering\arraybackslash 0,01 & \centering\arraybackslash 100 \\
\hline

\end{longtable}
\normalsize

\end{enumerate}


%%%%%%%%%%%%%%%%%%%%%%%%%%%%%%%%%%%%%%%%%%%%%%%%%% Цифровое выравнивание токов %%%%%%%%%%%%%%%%%%%%%%%%%%%%%%%%%%%%%%%%%%%%%%%%%%
\item Цифровое выравнивание токов

\begin{enumerate}[label=\arabic{section}.\arabic{subsection}.\arabic{enumi}.\arabic*, labelsep=4pt, leftmargin=0em, itemindent=65pt, parsep=0pt]

\item
В устройстве реализовано цифровое выравнивание токов присоединений, что позволяет подключать защиту к цепям ТТ с различными коэффициентами трансформации без использования промежуточных преобразователей.

\item
Токи присоединений приводятся к общему базисному току по выражению:
\begin{align}
i_n^{'}=\frac{i_n}{I_\textit{баз}}
\end{align}
\begin{eqexpl} [2mm]
    \item{$ i_n^{'} $} мгновенное значение выровненного тока n-го присоединения, о.е.;
	\item{$ i_n $} мгновенное значение первичного тока n-го присоединения, А;
	\item{$ I_\textit{баз} $} базисный ток, А.
\end{eqexpl}

\item
Величина базисного тока задается уставкой \textit{«Базисный ток»}. В качестве базисного тока выбирается наибольший номинальный первичный ток ТТ из всех присоединений:
\begin{align}
I_\textit{баз}=max(I_\textit{перв.ном,1}...I_\textit{перв.ном,n})
\end{align}

\end{enumerate}


%%%%%%%%%%%%%%%%%%%%%%%%%%%%%%%%%%%%%%%%%%%%%%%%%% Измерительные органы ДЗШ %%%%%%%%%%%%%%%%%%%%%%%%%%%%%%%%%%%%%%%%%%%%%%%%%%
\item Измерительные органы ДЗШ

\begin{enumerate}[label=\arabic{section}.\arabic{subsection}.\arabic{enumi}.\arabic*, labelsep=4pt, leftmargin=0em, itemindent=65pt, parsep=0pt]

\item
Каждый дифференциальный орган имеет идентичную реализацию и состоит из следующих измерительных органов:
\begin{itemize} 
\item дифференциальная токовая отсечка (ДТО);
\item дифференциальный орган с торможением (ДОТХ);
\item быстродействующий дифференциальный орган (БДО);
\item медленнодействующий дифференциальный орган (МДО);
\item блокировка при внешних коротких замыканиях (БВКЗ);
\item чувствительный токовый орган (ЧТО);
\item орган блокировки по второй гармонике дифференциального тока;
\item орган дифференциального тока для КЦТ.
\end{itemize}

\medskip

\begin{figure}[H]
\centering
\includegraphics[width=0.8\textwidth,height=0.8\textheight,keepaspectratio]{\fbpath/661.0631 ДЗШ/img/IO_DZSH.pdf}
\captionof{figure}{Функциональный блок ДЗШ ИО1}\label{pict:IO_DZSH}
\end{figure}

\medskip

\end{enumerate}


%%%%%%%%%%%%%%%%%%%%%%%%%%%%%%%%%%%%%%%%%%%%%%%%%% Дифференциальная токовая отсечка %%%%%%%%%%%%%%%%%%%%%%%%%%%%%%%%%%%%%%%%%%%%%%%%%%
\item Дифференциальная токовая отсечка

\begin{enumerate}[label=\arabic{section}.\arabic{subsection}.\arabic{enumi}.\arabic*, labelsep=4pt, leftmargin=0em, itemindent=65pt, parsep=0pt]

\item
ДТО предназначена для быстрого отключения защищаемого объекта при внутренних повреждениях, сопровождающихся большими токами.
Ввод в работу данной функции для всех дифференциальных органов осуществляется с помощью одного программного ключа \textbf{«ДТО»}.

\item
Данный ИО выполнен без торможения и может быть заблокирован только при обнаружении неисправности в ЦТ.
В качестве рабочей величины использует действующее значение первой гармоники дифференциального тока, которое определяется как модуль геометрической суммы первых гармоник токов присоединений, входящих в зону действия защиты:
\begin{align}
I_\textit{Д}=\mid \sum_{n=1}^N \vec{I}_{n}^{\,'} \mid
\end{align}
\begin{eqexpl} [2mm]
    \item{$ I_\textit{Д} $} действующее значение первой гармоники дифференциального тока, о.е.;
	\item{$ N $} количество присоединений, зафиксированных за соответствующим дифференциальным органом;
	\item{$ \vec{I}_{n}^{\,'} $} вектор первой гармоники выровненного тока n-го присоединения, о.е..
\end{eqexpl}

\item
Ток срабатывания ДТО для каждого дифференциального органа задается уставкой \textbf{«Iдто ПО»} (\textbf{«Iдто ИО1», «Iдто ИО2»}).


\end{enumerate}


%%%%%%%%%%%%%%%%%%%%%%%%%%%%%%%%%%%%%%%%%%%%%%%%%% Дифференциальный орган с торможением %%%%%%%%%%%%%%%%%%%%%%%%%%%%%%%%%%%%%%%%%%%%%%%%%%
\item Дифференциальный орган с торможением

\begin{enumerate}[label=\arabic{section}.\arabic{subsection}.\arabic{enumi}.\arabic*, labelsep=4pt, leftmargin=0em, itemindent=65pt, parsep=0pt]

\item
ДОТХ является основным измерительным органом и предназначен для отключения защищаемого объекта при всех видах внутренних повреждений.

\item
В качестве рабочих величин данный ИО использует действующие значения дифференциального и тормозного токов первой гармоники.
Действующее значение тормозного тока первой гармоники рассчитывается как полусумма модулей токов присоединений, входящих в зону действия защиты:
\begin{align}
I_\textit{Т}=0,5 \cdot \sum_{n=1}^N \mid \vec{I}_{n}^{\,'} \mid
\end{align}

\item
Тормозная характеристика дифференциальной защиты имеет два участка.
Горизонтальный участок не зависит от тормозного тока и необходим для отстройки от небольших токов небаланса при протекании малых сквозных токов, которые не вызывают насыщение ТТ.
Уровень срабатывания данного участка для каждого дифференциального органа задается отдельной уставкой \textbf{«Iд нач. ПО»} (\textbf{«Iд нач. ИО1»}, \textbf{«Iд нач. ИО2»}).
Наклонный участок обеспечивает отстройку от токов небаланса при насыщении ТТ.
Наклон характеристики определяется коэффициентом торможения, а точка начала наклона – значением начального тормозного тока.
Для каждого дифференциального органа эти параметры задаются отдельными уставками: \textbf{«Kт ПО»} и \textbf{«Iт нач. ПО»} (\textbf{«Kт ИО1»} и \textbf{«Iт нач. ИО1»}, \textbf{«Kт ИО2»} и \textbf{«Iт нач. ИО2»}).

\medskip

\begin{figure}[H]
\centering
\includegraphics[width=0.55\textwidth,height=0.55\textheight,keepaspectratio]{\fbpath/661.0631 ДЗШ/img/DiffZ_torm_char.pdf}
\captionof{figure}{Тормозная характеристика дифференциальной защиты}\label{pic:DiffZ_torm_char}
\end{figure}

\medskip

\end{enumerate}


%%%%%%%%%%%%%%%%%%%%%%%%%%%%%%%%%%%%%%%%%%%%%%%%%% Быстродействующий орган %%%%%%%%%%%%%%%%%%%%%%%%%%%%%%%%%%%%%%%%%%%%%%%%%%
\item Быстродействующий орган

\begin{enumerate}[label=\arabic{section}.\arabic{subsection}.\arabic{enumi}.\arabic*, labelsep=4pt, leftmargin=0em, itemindent=65pt, parsep=0pt]

\item
БДО предназначен для быстрого выявления внутренних повреждений до наступления насыщения ТТ.
Ввод в работу данной функции для всех дифференциальных органов осуществляется с помощью одного программного ключа \textbf{«БДО»}.

\item
Принцип действия органа основан на различии скорости нарастания дифференциального и тормозного токов при внутренних и внешних повреждениях.
При внутреннем КЗ обеспечивается равный рост дифференциального и тормозного тока в промежутке правильной трансформации ТТ.
При внешнем КЗ нарастание тормозного тока происходит раньше, тогда как рост дифференциального тока обусловлен насыщением ТТ.

\item
В алгоритме БДО реализована логика выявления одновременного нарастания фронтов действующих значений дифференциального и тормозного токов.
Сигнал срабатывания БДО формируется при наличии превышении дифференциального тока уставки срабатывания ДОТХ и нарастанием фронтов дифференциального и тормозного токов в пределах заданного времени.

\end{enumerate}


%%%%%%%%%%%%%%%%%%%%%%%%%%%%%%%%%%%%%%%%%%%%%%%%%% Медленнодействующий орган %%%%%%%%%%%%%%%%%%%%%%%%%%%%%%%%%%%%%%%%%%%%%%%%%%
\item Медленнодействующий орган

\begin{enumerate}[label=\arabic{section}.\arabic{subsection}.\arabic{enumi}.\arabic*, labelsep=4pt, leftmargin=0em, itemindent=65pt, parsep=0pt]

\item
МДО предназначен для выявления внутренних повреждений, а также для распознавания перехода КЗ из внешнего во внутреннее, в том числе в условиях насыщения ТТ.
Ввод в работу МДО осуществляется программным ключом \textbf{«МДО»}.

\item
Принцип действия органа основан на контроле формы абсолютного (выпрямленного) мгновенного дифференциального тока.
При внутреннем КЗ мгновенный дифференциальный ток изменяется два раза за период промышленной частоты, в то время как при внешнем КЗ с насыщением ТТ – один раз за период.
МДО обеспечивает требование по быстродействию при переходе внешнего КЗ во внутреннее.

\item
Выпрямленный мгновенный дифференциальный ток рассчитывается как модуль суммы мгновенных значений токов присоединений, входящих в зону действия защиты:
\begin{align}
i_\textit{д}=\mid \sum_{n=1}^N i_{n} \mid,
\end{align}
\begin{eqexpl} [0mm]
    \item{$ i_\textit{д} $} мгновенное значение выпрямленного дифференциального тока, о.е..
\end{eqexpl}

\end{enumerate}


%%%%%%%%%%%%%%%%%%%%%%%%%%%%%%%%%%%%%%%%%%%%%%%%%% Орган блокировки при внешнем КЗ %%%%%%%%%%%%%%%%%%%%%%%%%%%%%%%%%%%%%%%%%%%%%%%%%%
\item Орган блокировки при внешнем КЗ

\begin{enumerate}[label=\arabic{section}.\arabic{subsection}.\arabic{enumi}.\arabic*, labelsep=4pt, leftmargin=0em, itemindent=65pt, parsep=0pt]

\item
БВКЗ предназначена для выявления внешнего КЗ и обеспечения времени достоверного измерения тока на уровне 5 мс.
Ввод в работу данной функции для всех дифференциальных органов осуществляется с помощью одного программного ключа \textbf{«БВКЗ»}.

\item
Принцип действия органа основан на сравнении направлений токов присоединений.
Для идентификации внешнего повреждения формируются два вектора: вектор тока повреждённого присоединения ($ \vec{I}_1 $), вектор тока, рассчитываемого как разность между дифференциальным током и током повреждения ($ \vec{I}_2 $), т.е. вектор суммы токов всех остальных присоединений.

\item
В качестве вектора тока поврежденного присоединения выступает комплексное действующее значение основной гармоники тока того присоединения, в котором зафиксировано максимальное значение тока ($ \vec{I}_\textit{макс} $).

\begin{align}
\begin{split}
\vec{I}_1&=\vec{I}_\textit{макс}\\
\vec{I}_2&=\sum_{n=1}^N \vec{I}_{n} - \vec{I}_\textit{макс}
\end{split}
\end{align}

\item
Выбор присоединения с максимальным током основана на сравнении мгновенных значений токов всех присоединений.
В момент начала аварийного режима на выходе логики фиксируется комплексное действующее значение основной гармоники тока того присоединения, у которого на этот момент зафиксирован максимальный ток.
Данное присоединение определяется как поврежденное и его номер сохраняется на всё время действия КЗ.

\item
БВКЗ анализирует угол между двумя сформированными векторами.
При внешнем КЗ угол между векторами превышает 90º.
Данный признак используется как критерий для формирования сигнала блокировки.

\medskip

\begin{figure}[H]
\centering
\begin{subfigure}{0.45\textwidth}
    \centering
    \includegraphics[width=\textwidth, trim=0 0 0 0, clip]{\fbpath/661.0631 ДЗШ/img/BVKZ_a.pdf}
    \caption{а) внутреннее КЗ}
    \label{pict:Action_DFZ_a}
\end{subfigure}
%\hfill
\begin{subfigure}{0.45\textwidth}
    \centering
    \includegraphics[width=\textwidth, trim=0 0 0 0, clip]{\fbpath/661.0631 ДЗШ/img/BVKZ_b.pdf}
    \caption{б) внешнее КЗ}
    \label{pict:Action_DFZ_b}
\end{subfigure}

\caption{Характеристика срабатывания БВКЗ}
\label{pict:BVKZ}
\end{figure}

\medskip

\item
Действие БВКЗ блокируется в режиме опробования шин, поскольку в данном режиме контролируется ток только одного присоединения, что исключает возможность корректного сравнения направлений векторов.

\end{enumerate}


%%%%%%%%%%%%%%%%%%%%%%%%%%%%%%%%%%%%%%%%%%%%%%%%%% Чувствительный токовый орган %%%%%%%%%%%%%%%%%%%%%%%%%%%%%%%%%%%%%%%%%%%%%%%%%%
\item Чувствительный токовый орган

\begin{enumerate}[label=\arabic{section}.\arabic{subsection}.\arabic{enumi}.\arabic*, labelsep=4pt, leftmargin=0em, itemindent=65pt, parsep=0pt]

\item
В устройстве предусмотрено два подхода для надежного отключения выключателей присоединений при работе ДЗШ, в том числе в режимах АПВ шин и опробования.
Такие режимы могут сопровождаться пониженным уровнем токов КЗ, в результате чего чувствительность дифференциальных органов может быть недостаточной.
Первый подход подразумевает применение ЧТО. Второй подход заключается в очувствлении уставок тормозной характеристики, которые используются в ДОТХ, БДО и МДО.

\item
ЧТО выполнен в виде реле тока, включенного на дифференциальные токи пускового и избирательных органов.
Ток срабатывания ЧТО для каждого дифференциального органа задается отдельной уставкой \textbf{«Iчто ПО»} (\textbf{«Iчто ИО1»}, \textbf{«Iчто ИО2»}).

\item
Очувствление представляет собой автоматическую смену величины начального тока срабатывания тормозной характеристики на более чувствительную уставку.
Для каждого дифференциального органа параметры тормозной характеристики в режиме очувствления задаются отдельными уставками.
Подробное описание алгоритма очувствления приведен в разделе «Логика очувствления».

\end{enumerate}


%%%%%%%%%%%%%%%%%%%%%%%%%%%%%%%%%%%%%%%%%%%%%%%%%% Блокировка по второй гармонике дифференциального тока %%%%%%%%%%%%%%%%%%%%%%%%%%%%%%%%%%%%%%%%%%%%%%%%%%
\item Блокировка по второй гармонике дифференциального тока

\begin{enumerate}[label=\arabic{section}.\arabic{subsection}.\arabic{enumi}.\arabic*, labelsep=4pt, leftmargin=0em, itemindent=65pt, parsep=0pt]

\item
Для предотвращения срабатывания ИО ДЗШ (за исключением ДТО) при БНТ, возникающих при включении присоединения с силовым трансформатором, используется блокировка по второй гармонике дифференциального тока.
Влияние второй гармоники, возникающей при БНТ, на работу ИО ДЗШ проявляется только в условиях опробования по схеме с «открытым плечом».

\item
ИО блокировки выполнен в виде реле тока, реагирующее на отношение дифференциального тока второй гармоники к дифференциальному току первой гармоники.
Величина срабатывания для каждого дифференциального органа задается отдельной уставкой \textbf{«I2г/I1г ПО»} (\textbf{«I2г/I1г ИО1»}, \textbf{«I2г/I1г ИО2»}).

\item
При срабатывании ИО ДЗШ сигнал на отключение сформируется по истечению соответствующей выдержки времени \textbf{«Т ДЗШ ПО»} (\textbf{«Т ДЗШ ИО1»}, \textbf{«Т ДЗШ ИО2»}).
Уставка времени срабатывания задается отдельно для каждого дифференциального органа.

\item
При нарушении фиксации присоединений сигналы отключения от ДЗШ и ЧТО формируются для обеих СШ.

\medskip

\begin{figure}[H]
\centering
\includegraphics[width=0.9\textwidth,height=0.9\textheight,keepaspectratio]{\fbpath/661.0631 ДЗШ/img/DZSH.pdf}
\captionof{figure}{Функциональный блок ДЗШ}\label{pict:DZSH}
\end{figure}

\medskip

\end{enumerate}


%%%%%%%%%%%%%%%%%%%%%%%%%%%%%%%%%%%%%%%%%%%%%%%%%% Контроль цепей тока %%%%%%%%%%%%%%%%%%%%%%%%%%%%%%%%%%%%%%%%%%%%%%%%%%
\item Контроль цепей тока

\begin{enumerate}[label=\arabic{section}.\arabic{subsection}.\arabic{enumi}.\arabic*, labelsep=4pt, leftmargin=0em, itemindent=65pt, parsep=0pt]

\item
В устройстве реализован контроль исправности токовых цепей для исключения ложного срабатывания ДЗШ вследствие нарушения баланса токов в дифференциальной цепи при обрыве, замыканий и т.д.

\medskip

\begin{figure}[H]
\centering
\includegraphics[width=0.9\textwidth,height=0.9\textheight,keepaspectratio]{\fbpath/661.0631 ДЗШ/img/KCT.pdf}
\captionof{figure}{Функциональный блок КЦТ}\label{pict:KCT}
\end{figure}

\medskip

\item
Ввод и вывод контроля исправности токовых цепей для каждого дифференциального органа осуществляется отдельным программным ключом \textbf{«КЦТ ПО»} (\textbf{«КЦТ ИО1»}, \textbf{«КЦТ ИО2»}).
В режиме нарушенной фиксации КЦТ избирательных органов выводится из работы.

\item
Для выявления неисправности токовых цепей предусмотрено три измерительных органа, реагирующих на дифференциальные токи ПО, ИО1 и ИО2.
Ток срабатывания КЦТ для каждого дифференциального органа задается отдельной уставкой \textbf{«Iкцт ПО»} (\textbf{«Iкцт ИО1»}, \textbf{«Iкцт ИО2»}).

\item
При обнаружении небаланса в токовых цепях измерительный орган с соответствующей выдержкой времени \textbf{«Ткцт ПО»} (\textbf{«Ткцт ИО1»}, \textbf{«Ткцт ИО2»}) действует на сигнализацию, а также на блокировку ДЗШ.
Блокировка ДЗШ от КЦТ осуществляется при введенном программном ключе \textbf{«Реж. КЦТ ПО»} (\textbf{«Реж. КЦТ ИО1»}, \textbf{«Реж. КЦТ ИО2»}).
Данные программные ключи приведены в перечне уставок функции ДЗШ.

\item
Сброс блокировки выполняется внешним сигналом \textit{«Возврат КЦТ»}.
Оперативный вывод блокировки ДЗШ при неисправностях в токовых цепях осуществляется внешним сигналом \textbf{«Вывод блок. ДЗШ от КЦТ»}.


\small
\begin{longtable}{|>{\centering\arraybackslash}m{5.3cm}|>{\centering\arraybackslash}m{3.3cm}|>{\centering\arraybackslash}m{4.2cm}|>{\centering\arraybackslash}m{1.8cm}|>{\centering\arraybackslash}m{1cm}|}
% Первая страница: полный заголовок
\caption{Перечень уставок функции ДЗШ\hfill\vspace{-0.5\baselineskip}}\label{desc:tbl2}\\
\hline
\rowcolor{gray!30}
Параметр & Условное обозначение на схеме & Значение/ Диапазон & Единица измерения & Шаг \\
\hline
\endfirsthead

% Промежуточные страницы: «Продолжение таблицы»
\caption*{\hspace{3pt}\emph{Продолжение таблицы \ref{desc:tbl2}}\hfill\ }\vspace{-0.5\baselineskip} \\
\hline
\rowcolor{gray!30}
Параметр & Условное обозначение на схеме & Значение/ Диапазон & Единица измерения & Шаг \\
\hline
\endhead

%===m>ДЗШ|-
%===m
\end{longtable}
\normalsize	

\end{enumerate}

\end{enumerate}
